\section*{Erd\H{o}s Problem \#357: distinct consecutive sums}
\subsection*{1. Statement and scope}
Let $f(n)$ be the largest $k$ for which some
$1\le a_1<\cdots<a_k\le n$ has all sums
$\sum_{i=u}^v a_i$, $1\le u\le v\le k$, distinct.
How does $f(n)$ grow, and is $f(n)=o(n)$?
Related questions ask whether an infinite such sequence has density
zero and whether its reciprocal sum converges.
We prove explicit lower and upper bounds and an infinite-sequence
growth obstruction. The sublinear upper bound and reciprocal-sum
questions remain unresolved.

\subsection*{2. An explicit lower bound of order $2\sqrt n$}
Put $m=\lfloor\sqrt n\rfloor$, $k=2m-1$, and take the last $k$ integers
up to $n$, namely $a_i=n-k+i$. These are positive since
$m^2\ge2m-1$. Sums of a fixed length $r$ strictly increase with the
starting index. Their minimum and maximum are
\[
 L_r=r(n-k)+r(r+1)/2,\qquad U_r=rn-r(r-1)/2.
\]
For $1\le r<k$,
\[
 L_{r+1}-U_r
 =n-(r+1)(k-1)+r^2
 =(r-m+1)^2+n-m^2+1\ge1.
\]
Thus the ranges of different lengths are disjoint, proving
\[
 f(n)\ge2\lfloor\sqrt n\rfloor-1.
\]
This is the familiar interval construction also checked in the
earlier GPT Pro 5.2 attempt.

\subsection*{3. An elementary upper bound with controlled error}
Write $S=\sum_{i=1}^k a_i$. For $1\le h\le k$, there are
\[
 M=\sum_{r=1}^h(k-r+1)=hk-h(h-1)/2
\]
short consecutive sums of lengths at most $h$. They are distinct
positive integers, so their total is at least $M(M+1)/2$.
Each $a_i$ occurs in at most $r$ sums of length $r$, so the total is
at most $h(h+1)S/2$. Therefore
\[
 M^2\le h(h+1)S,\qquad
 (k-h/2)^2\le(1+1/h)S.                            \tag{1}
\]
Choose $h=\lfloor\sqrt k\rfloor$. If $S\ge k^2$ there is nothing
to prove; otherwise (1) implies
\[
 k^2\le S+kh+k^2/h=S+O(k^{3/2}).
\]
The individual $a_i$ are distinct and bounded by $n$, giving
$S\le kn-k(k-1)/2$. Combining these inequalities yields
\[
 f(n)\le\frac23 n+O(\sqrt n).                      \tag{2}
\]
The argument even permits a nonmonotone ordering, because distinct
length-one sums still force distinct entries. It reconstructs a
classical weaker upper bound; it is not asserted to be the best
published constant.

\subsection*{4. A quantitative obstruction for an infinite sequence}
For any infinite increasing sequence with the stated property,
\[
 \limsup_{j\to\infty}\frac{a_j}{j\log j}\ge\frac12.                  \tag{3}
\]
Suppose instead that $a_j\le c j\log j$ for all sufficiently large
$j$, with $c<1/2$. Choose a fixed small $\lambda>0$ so that
$1/(c(1+\lambda))>2$.
For large starting indices $j\le N/(1+\lambda)$ and
$1\le r\le\lfloor\lambda j\rfloor$, monotonicity gives
\[
 \sum_{i=j}^{j+r-1}a_i
 \le c(1+\lambda)rj\log((1+\lambda)j).
\]
All these sums are distinct positive integers at most
$S_N=a_1+\cdots+a_N=O(N^2\log N)$. Their reciprocal sum is therefore
at most the harmonic sum $H_{S_N}=(2+o(1))\log N$. On the other hand
it is at least
\[
 \frac1{c(1+\lambda)}
 \sum_{j_0\le j\le N/(1+\lambda)}
 \frac{H_{\lfloor\lambda j\rfloor}}{j\log((1+\lambda)j)}
 =\left(\frac1{c(1+\lambda)}+o(1)\right)\log N.
\]
The last equality follows from
$H_{\lfloor\lambda j\rfloor}=\log j+O_\lambda(1)$; the accumulated
error is $O_\lambda(\log\log N)$.
The two bounds contradict the choice of $\lambda$, proving (3).
In particular $j/a_j$ tends to zero along a subsequence, so the
set of terms has lower asymptotic density zero.

\subsection*{5. Assessment and references}
The interval construction and (2) leave a large gap between
$\sqrt n$ and $n$. Statement (3) concerns a subsequence in one fixed
infinite example and does not yield a uniform bound for all finite
examples. Lower density zero is not density zero; nor does (3)
imply convergence of $\sum1/a_j$.

Source: \texttt{https://www.erdosproblems.com/357}, checked 2026-09-25,
which credits Erd\H{o}s--Harzheim for the question and records the
averaging obstruction and Hegyv\'ari's bounds.
The prior GPT Pro 5.2 constructions were checked and retained with
explicit attribution. The arguments here use no unproved external
theorem, and no novelty or formal Lean verification is claimed.
\subsection*{COMPLETION ESTIMATE}
\textbf{COMPLETION: 50\%}
